%% kapa: exam=1 round=37 year=2026 period=1 form=A subject=civil_law subject_ko=민법 seq=27 number=27
%% source: Q-Net 감정평가사 기출문제 (cst00309) artlSeq=5258924 file=raw/1008/5258924/2026년 제37회 감정평가사 1차 1교시.pdf
%% conversion: kordoc 4.24.51 -> tools/gen1.py (auto); review=proofread-sonnet; flags=none
\begin{question}{27}
점유자와 회복자의 관계에 관한 설명으로 옳은 것을 모두 고른 것은? (다툼이 있으면 판례에 따름)
\begin{framedtext}
ㄱ. 선의의 점유자는 비록 법률상 원인 없이 타인의 건물을 점유ㆍ사용하고 이로 말미암아 그에게 손해를 입혔다고 하더라도 그 점유ㆍ사용으로 인한 이득을 반환할 의무는 없다. \par
ㄴ. 농지매매계약이 무효가 되어 그 농지의 선의의 매수인이 매도인에게 농지를 반환하는 경우, 매수인이 그 농지를 경작하여 수확을 하였다면 매수인은 그 농지에 투입한 통상필요비인 비료비용을 매도인에게 상환청구할 수 없다. \par
ㄷ. 점유자가 점유물의 현상을 유지하기 위하여 지출한 금액에 대하여 법원은 회복자의 청구에 의하여 상당한 상환기간을 허여하여야 한다. \par
\end{framedtext}
\begin{choices}
  \item ㄱ
  \item ㄴ
  \item ㄱ, ㄴ
  \item ㄴ, ㄷ
  \item ㄱ, ㄴ, ㄷ
\end{choices}
\end{question}
\answer{3}
